\chapter{Homologia i kohomologia \v{C}echa: od pokryć do granic przestrzeni}
\label{ch:cech-v4}

\section{Motywacja: co dzieje się przy nieskończonym przybliżaniu?}
Homologia singularna buduje łańcuchy z pojedynczych ciągłych map
standardowych sympleksów do $X$. Dobrze współpracuje z homotopią,
ale w przestrzeni o coraz drobniejszej, nieskończonej strukturze
lokalnej nie zawsze zachowuje się jak granica homologii
skończonych przybliżeń. Można zamiast tego obserwować $X$
przez \emph{otwarte pokrycia}: każdy układ przecinających się
zbiorów opisujemy skończonym kompleksem symplicjalnym,
a następnie pokrycia sukcesywnie zagęszczamy.

Ta zmiana punktu widzenia rozwiązuje szczególnie ważny problem:
jak obliczyć niezmienniki zwartego zbioru $A$ za pomocą malejących
otoczeń $U\supset A$ i jak związać je z homologią dopełnienia
$S^n\setminus A$ (dualność Aleksandra). Trzeba jednak od razu
rozróżnić \emph{homologię} i \emph{kohomologię} \v{C}echa:
pierwsza bierze granicę odwrotną, druga granicę prostą.

\section{Nerw pokrycia i odwzorowania przy rozdrobnieniu}
\begin{definicja}[Nerw]
Niech $\mathcal U=(U_i)_{i\in J}$ będzie pokryciem otwartym
przestrzeni $X$. \emph{Nerwem} $N(\mathcal U)$ jest abstrakcyjny
kompleks symplicjalny o wierzchołkach $i\in J$; skończona rodzina
różnych wierzchołków $\{i_0,\ldots,i_k\}$ wyznacza sympleks
wtedy i tylko wtedy, gdy $U_{i_0}\cap\cdots\cap U_{i_k}\ne\varnothing$.
Do skończonego pokrycia przyporządkowujemy zwykły skończony
kompleks symplicjalny, do którego stosujemy homologię
symplicjalną oraz kohomologię.
\end{definicja}

\begin{przyklad}[Okrąg]
Pokryjmy $S^1$ trzema dostatecznie krótkimi otwartymi łukami
$U_0,U_1,U_2$ tak, aby każde dwa się przecinały, lecz
$U_0\cap U_1\cap U_2=\varnothing$.
Nerw ma trzy wierzchołki i trzy krawędzie, ale nie ma
wypełnionego trójkąta; jego realizacja geometryczna jest więc
homeomorficzna z $S^1$.
Jeżeli dopuścimy niepuste przecięcie potrójne, pojawi się
$2$-sympleks, a nerw stanie się ściągalny. Samo pojedyncze
pokrycie nie może zatem definiować niezmiennika topologicznego:
trzeba rozważać \emph{wszystkie} coraz drobniejsze pokrycia.
\end{przyklad}

\begin{definicja}[Rozdrobnienie]
Pokrycie $\mathcal V$ jest \emph{drobniejsze} od $\mathcal U$,
co zapisujemy $\mathcal V\succ\mathcal U$, jeśli dla każdego
$V\in\mathcal V$ istnieje $U\in\mathcal U$ z $V\subseteq U$.
Wybór takich przyporządkowań daje mapę symplicjalną
$r_{\mathcal V\mathcal U}:N(\mathcal V)\to N(\mathcal U)$:
każdy wierzchołek $V$ wysyłamy na wybrany $U$ zawierający $V$.
\end{definicja}

\begin{lemat}[Niezależność wyborów]
\label{lem:cech-contiguity-v4}
Dwie mapy nerwów wybrane dla jednego rozdrobnienia
$\mathcal V\succ\mathcal U$ są przyległe symplicjalnie,
a zatem indukują te same homomorfizmy homologii i kohomologii.
Przy trzykrotnym rozdrobnieniu $\mathcal W\succ\mathcal V\succ
\mathcal U$ mapy indukowane są zgodne ze składaniem.
\end{lemat}
\begin{proof}
Niech $\{V_0,\ldots,V_k\}$ będzie sympleksem nerwu
$N(\mathcal V)$ i $x\in\bigcap_j V_j$.
Jeśli $U_j$ i $U'_j$ to dwa możliwe wybory zbiorów z
$\mathcal U$ zawierających $V_j$, to wszystkie $U_j,U'_j$
zawierają $x$. Ich wierzchołki tworzą więc wspólnie sympleks
w $N(\mathcal U)$, co jest dokładnie warunkiem przyległości.
Homotopia liniowa między realizacjami map przyległych
pozostaje w tym samym sympleksie, jest więc ciągła.
Mapy homotopijne działają identycznie na homologię i kohomologię.
Przy składaniu oba wybory z $\mathcal W$ do $\mathcal U$
również są przyległe, stąd ostatnie twierdzenie.
\end{proof}

\section{Definicje i kierunki granic}
W dalszym ciągu przyjmujemy, że $X$ jest \emph{zwartą przestrzenią
Hausdorffa}, a $G$ grupą abelową. Wszystkie otwarte pokrycia
mają skończone podpokrycia, a każde dwa skończone pokrycia
mają wspólne drobniejsze pokrycie (ich niepuste przecięcia).
Zatem skończone pokrycia tworzą zbiór skierowany przez
relację rozdrobnienia.

\begin{definicja}[Homologia i kohomologia \v{C}echa]
Dla każdego $n\ge0$ definiujemy
\[
\check H_n(X;G)=\varprojlim_{\mathcal U} H_n(N(\mathcal U);G),
\qquad
\check H^n(X;G)=\varinjlim_{\mathcal U}H^n(N(\mathcal U);G).
\]
W pierwszym układzie mapy $H_n(N(\mathcal V);G)\to
H_n(N(\mathcal U);G)$ biegną od drobniejszego do grubszego
pokrycia. W drugim układzie przez kontrawariantność
kohomologii otrzymujemy mapy
$H^n(N(\mathcal U);G)\to H^n(N(\mathcal V);G)$.
\emph{Zredukowane} grupy $\widetilde{\check H}_n$ i
$\widetilde{\check H}^{\,n}$ definiuje się za pomocą
zredukowanych (ko)homologii nerwów.
\end{definicja}

\begin{uwaga}[Co dosłownie znaczy granica odwrotna?]
Elementem $\check H_n(X;G)$ jest rodzina
$(a_{\mathcal U})_{\mathcal U}$, gdzie
$a_{\mathcal U}\in H_n(N(\mathcal U);G)$, spełniająca
$(r_{\mathcal V\mathcal U})_*a_{\mathcal V}=a_{\mathcal U}$
przy $\mathcal V\succ\mathcal U$.
Nie wystarcza wybrać dowolny cykl dla każdego pokrycia:
klasy muszą być zgodne przy rozdrobnieniu.
Natomiast element kohomologii \v{C}echa reprezentuje pewna
klasa w nerwie \emph{jednego} pokrycia; klasy reprezentują
ten sam element, jeśli po przejściu do wspólnego drobniejszego
pokrycia stają się równe. Jest to istota granicy prostej.
\end{uwaga}

\begin{stwierdzenie}[Funktorialność]
Jeśli $f:X\to Y$ jest ciągła między zwartymi przestrzeniami
Hausdorffa, to indukuje
$f_*:\check H_n(X;G)\to\check H_n(Y;G)$
oraz $f^*:\check H^n(Y;G)\to\check H^n(X;G)$.
Odwzorowania te są zgodne ze składaniem i identycznością.
\end{stwierdzenie}
\begin{proof}
Dla skończonego pokrycia $\mathcal U$ przestrzeni $Y$
przeciwobrazy $f^{-1}(U)$ stanowią skończone pokrycie $X$.
Zachodzi naturalna mapa nerwów
$N(f^{-1}\mathcal U)\to N(\mathcal U)$ (puste przeciwobrazy
pomijamy). Zgodność z rozdrobnieniami wynika z
Lematu~\ref{lem:cech-contiguity-v4}. Przechodząc do granicy
odwrotnej, otrzymujemy $f_*$; po odwróceniu strzałek
kohomologii dostajemy $f^*$.
\end{proof}

\section{Porównanie z homologią singularną}
\begin{stwierdzenie}[Porównanie dla wielościanów]
\label{prop:cech-poly-v4}
Jeżeli $X$ jest zwartym wielościanem (w szczególności
skończonym CW-kompleksem, który jest triangulowalny), to
naturalne homomorfizmy porównawcze są izomorfizmami
$H_n(X;G)\cong\check H_n(X;G)$ i
$H^n(X;G)\cong\check H^n(X;G)$ dla $n\ge0$.
\end{stwierdzenie}
\begin{proof}
Ustalmy skończoną triangulację $K$. Pokrycia otwartymi
gwiazdami w kolejnych podziałach barycentrycznych są
kofinalne wśród skończonych pokryć otwartych: z lematu
Lebesgue'a o liczbie pokrycia i malejących średnic sympleksów
wynika, że po dostatecznej liczbie podziałów każda taka
gwiazda mieści się w jednym zbiorze zadanego pokrycia.
Nerw pokrycia otwartymi gwiazdami wierzchołków odpowiedniego
podziału jest izomorficzny z jego kompleksem symplicjalnym.
Mapy nerwów dla kolejnych rozdrobnienień odpowiadają
mapom aproksymacji symplicjalnej identyczności $|K|\to|K|$,
a zatem indukują izomorfizmy na (ko)homologii.
Układy odwrotny i prosty stają się więc układami
stałymi z dokładnością do kanonicznych izomorfizmów.
\end{proof}

\begin{uwaga}
Nie można bez dodatkowych hipotez ogłosić powyższych
izomorfizmów dla wszystkich przestrzeni zwartych.
W szczególności spójność i spójność łukowa nie są tym samym:
$H_0^{\mathrm{sing}}(X;\Z)$ wykrywa składowe łukowe,
natomiast $\check H_0$ w zwartych przykładach z poniższej
sekcji może rejestrować jedynie spójność widoczną w
kolejnych skończonych przybliżeniach.
\end{uwaga}

\section{Kolczyk hawajski: coraz więcej okręgów}
\begin{definicja}
Kolczykiem hawajskim nazywamy zwarty podzbiór płaszczyzny
$\mathbb H=\bigcup_{n\ge1}C_n$, gdzie
$C_n=\{(x,y)\in\R^2:(x-1/n)^2+y^2=1/n^2\}$.
Wszystkie okręgi są styczne w $(0,0)$, a ich średnice
maleją do zera.
\end{definicja}

\begin{stwierdzenie}
Dla każdego $k\ge1$ zachodzi
$\check H_1(\mathbb H;\Z)\cong\prod_{n=1}^{\infty}\Z$,
zaś $\check H_k(\mathbb H;\Z)=0$ dla $k\ge2$.
\end{stwierdzenie}
\begin{proof}
Niech $X_m=C_1\vee\cdots\vee C_m$ będzie bukietem
pierwszych $m$ okręgów. Zbiory coraz mniejszych ogonów
$\bigcup_{n>m}C_n$ kurczą się w dowolnie małym otoczeniu
punktu wspólnego. Pozwala to dobrać kofinalny układ
skończonych pokryć $\mathbb H$, których nerwy mają typ
homotopijny $X_m$, a mapy od etapu $m+1$ do $m$ odpowiadają
zgnieceniu ostatniego okręgu. W szczególności układ
pierwszych grup homologii to
\[
\Z\longleftarrow\Z^2\longleftarrow\Z^3\longleftarrow\cdots,
\qquad (a_1,\ldots,a_{m+1})\longmapsto(a_1,\ldots,a_m).
\]
Rodzina zgodna jest wyznaczona przez nieskończony ciąg
liczb całkowitych $(a_1,a_2,\ldots)$, zatem granica odwrotna
to iloczyn $\prod_{n\ge1}\Z$. Dla $k\ge2$ wszystkie grupy
w tym układzie są zerowe. Uzasadnienie kofinalności można
przeprowadzić, dobierając dla danego pokrycia otoczenie
punktu $(0,0)$ mieszczące wszystkie dostatecznie małe okręgi
i dzieląc skończoną liczbę pozostałych łuków.
\end{proof}

\begin{uwaga}[Nie mylić z grupą fundamentalną]
Nie stwierdzamy, że $\pi_1(\mathbb H)$ jest wolną grupą o
przeliczalnej liczbie generatorów. Występują pętle
przebiegające nieskończenie wiele okręgów w coraz krótszym
czasie. To właśnie typ zjawiska, przy którym model przez
układ skończonych nerwów daje inną perspektywę niż
pojedynczy skończony kompleks CW.
\end{uwaga}

\section{Solenoid: granica okręgów i zanik pierwszej homologii}
Niech $p\ge2$ będzie liczbą całkowitą. Definiujemy
\emph{solenoid $p$-adyczny}
$\Sigma_p=\varprojlim(S^1\xleftarrow{z\mapsto z^p}S^1
\xleftarrow{z\mapsto z^p}\cdots)$.
Punktami solenoidu są ciągi $(z_0,z_1,\ldots)$ spełniające
$z_j=z_{j+1}^p$. Jest to zwarta przestrzeń Hausdorffa,
a każdy rzut na skończony poziom jest ciągły.

\begin{twierdzenie}[Dwie granice o przeciwnych wynikach]
\label{thm:solenoid-v4}
Dla $\Sigma_p$ zachodzi
\[
\check H_0(\Sigma_p;\Z)\cong\Z,\qquad
\check H_1(\Sigma_p;\Z)=0,\qquad
\check H^{1}(\Sigma_p;\Z)\cong\Z[1/p].
\]
Solenoid jest spójny, ale nie jest łukowo spójny.
\end{twierdzenie}
\begin{proof}
W przypadku granic odwrotnych zwartych wielościanów
$X=\varprojlim X_j$ z surjektywnymi mapami przejścia
pokrycia wstecznymi obrazami pokryć poziomów tworzą
rodzinę kofinalną. Wynika to z faktu, że ich obrazy zwrotne
tworzą bazę topologii granicy, a zwartość pozwala przejść
od pokryć do skończonych podpokryć. Stąd
$\check H_n(X;G)\cong\varprojlim H_n(X_j;G)$ oraz
$\check H^n(X;G)\cong\varinjlim H^n(X_j;G)$
dla tego układu wielościanów.

Na każdym poziomie mamy $H_0(S^1;\Z)=\Z$ i
$H_1(S^1;\Z)=\Z$. Odwzorowanie $z\mapsto z^p$ działa
na $H_0$ identycznie, a na $H_1$ mnożeniem przez $p$.
Zatem element $\check H_1$ to rodzina
$(a_0,a_1,\ldots)$ z $a_j=pa_{j+1}$.
Wymaga to, aby $a_0$ było podzielne przez dowolnie
wysoką potęgę $p$, zatem wszystkie $a_j=0$.
W kohomologii strzałki biegną odwrotnie, więc
$\check H^1$ jest granicą prostą
$\Z\xrightarrow{\cdot p}\Z\xrightarrow{\cdot p}\cdots$.
Klasa liczby $m$ z poziomu $j$ odpowiada $m/p^j$,
co daje $\Z[1/p]$.

Spójność wynika z tego, że granica odwrotna niepustych
zwartych przestrzeni spójnych przy surjektywnych mapach
przejścia jest spójna: gdyby granica rozpadła się na dwa
niepuste domknięto-otwarte zbiory, zwartość pozwoliłaby
wykryć to rozdzielenie na dostatecznie wysokim poziomie,
wbrew spójności okręgu.

Dla braku łukowej spójności rozważmy włókno rzutu
$\Sigma_p\to S^1$ nad $1$. Jest ono homeomorficzne z
$\Z_p$, zbiorem ciągów reszt zgodnych modulo $p^j$.
Ścieżka od punktu $(1,1,\ldots)$ do punktu tego włókna
po złożeniu z pierwszym rzutem daje pętlę o całkowitej
liczbie owinięć $m$. Kolejne współrzędne ścieżki są
jednoznacznymi podniesieniami tej pętli dla pokryć stopnia
$p^j$, zatem jej końcowe współrzędne odpowiadają resztom
$m\bmod p^j$. Możliwych takich końców jest przeliczalnie
wiele (po jednym dla $m\in\Z$), podczas gdy $\Z_p$ jest
nieprzeliczalny. Stąd nie każde dwa punkty można połączyć
ścieżką.
\end{proof}

\begin{wniosek}
Dla solenoidu $H_0^{\mathrm{sing}}(\Sigma_p;\Z)$ nie jest
izomorficzna z $\check H_0(\Sigma_p;\Z)=\Z$:
pierwsza grupa jest wolną grupą na składowych łukowych,
których jest więcej niż jedna. Nie należy zatem bez
hipotez lokalnej regularności zamieniać teorii singularnej
na teorię \v{C}echa.
\end{wniosek}

\section{Kohomologia \v{C}echa przez otoczenia i rola dualności}
\begin{twierdzenie}[Opis przez otoczenia w sferze]
\label{thm:cech-neighbourhood-v4}
Niech $A\subset S^n$ będzie niepustym zbiorem zwartym,
$G$ grupą abelową. Wówczas
$\check H^k(A;G)\cong\varinjlim_{U\supset A}H^k(U;G)$,
gdzie granicę bierzemy po otwartych otoczeniach $U$ zbioru
$A$, przy przejściach danych przez restrykcje.
\end{twierdzenie}
\begin{proof}
Wybieramy triangulacje sfery o coraz mniejszych średnicach
komórek. Dla ustalonego otwartego $U\supset A$ zwartość $A$
pozwala dobrać skończone pokrycie otwartymi gwiazdami
w dostatecznie drobnym podziale, którego suma leży w $U$.
Po niewielkim pogrubieniu otrzymujemy regularne otoczenie
$N$ zbioru $A$ o typie homotopijnym nerwu tego pokrycia
(lemat o nerwie, ponieważ wszystkie niepuste przecięcia
otwartych gwiazd są ściągalne). Dla odwrotnego przejścia,
każde skończone pokrycie otwarte $A$ można przedłużyć do
pokrycia pewnego otoczenia $U$ poprzez wybór otwartych
zbiorów sfery mających zadane ślady na $A$; po zwężeniu
otoczenia nie pojawiają się nowe niepuste przecięcia
skończonych rodzin, które były puste na $A$.
Te dwa kroki dają kofinalne układy nerwów i otoczeń,
a więc identyczną granicę prostą kohomologii.
\end{proof}

\begin{uwaga}[Dlaczego nie pojawia się tu homologia singularna $A$?]
W dualności Aleksandra dla dowolnego zwartego $A\subset S^n$
po prawej stronie musi stać kohomologia \v{C}echa.
Dla skończonych wielościanów oba rodzaje kohomologii
się zgadzają (Stwierdzenie~\ref{prop:cech-poly-v4}),
lecz w przypadku zbiorów dzikich nie wolno wykonywać
takiego podstawienia. Granica \emph{prosta} w powyższym
wzorze jest również zgodna z kierunkiem map wywołanym
przez malejące otoczenia.
\end{uwaga}

\section{Granice odwrotne a dokładność}
Granica prosta grup abelowych jest funktorem dokładnym,
ale granica odwrotna w ogólności nie jest prawostronnie
dokładna. Wynika stąd, że naiwnie zdefiniowana homologia
\v{C}echa nie musi spełniać wszystkich aksjomatów dokładności
homologii Eilenberga--Steenroda dla dowolnych par zwartych
przestrzeni. Przy bardziej zaawansowanym badaniu granic
stosuje się pochodny funktor $\varprojlim{}^1$ i tzw.
silne homologie. W tej książce rozróżniamy te konstrukcje,
nie utożsamiając ich bez dowodu.

\begin{uwaga}
Dla zadanego układu odwrotnego grup abelowych z
surjektywnymi mapami przejścia (przy przeliczalnym
indeksowaniu) zachodzą przydatne warunki zaniku
$\varprojlim{}^1$. Nie należy jednak przypisywać tej
własności każdemu układowi pochodzącemu od pokryć;
przykład solenoidu ma na $H_1$ mapy mnożenia przez $p$,
które nie są surjekcjami.
\end{uwaga}
\clearpage

